\section{A Taxonomy of Drift (RQ1)}
\label{sec:taxonomy}

\begin{figure}[t]
\centering
\resizebox{\columnwidth}{!}{% Approved intent vs runtime reality (single column, resized)
\begin{tikzpicture}[node distance=3.2mm and 4mm]
  % approval event
  \node[human, minimum width=20mm] (appr) at (0,0)
    {approval event $t_0$\\{\smlbl snapshot $\Gapp(t_0)$}};
  \node[attn, below=2.5mm of appr, minimum width=20mm] (snap)
    {$m_0$, $\delta_0$, $\dig(\pi_0)$, $A_0$, $\iota_0$, $\sigma_0$};

  % runtime plane
  \node[stage, right=8mm of appr, minimum width=17mm] (r1)
    {observed\\{\smlbl config $m_t$}};
  \node[stage, right=3.5mm of r1, minimum width=15mm] (r2)
    {running\\{\smlbl digests $\delta_t$}};
  \node[stage, right=3.5mm of r2, minimum width=15mm] (r3)
    {environment\\{\smlbl $\varepsilon_t$}};
  \node[lbl, above=1mm of r2] {\textbf{runtime plane}};
  \draw[flow] (r1) to[out=-25,in=205] node[lbl, below, pos=0.22]{reconciler} (r3);
  \draw[flow] (r3) to[out=155,in=25] (r1);

  % governance plane
  \node[evid, below=11mm of r1, minimum width=17mm] (g1)
    {policy\\{\smlbl $\Pol(t)$}};
  \node[evid, below=11mm of r2, minimum width=15mm] (g2)
    {authorization\\{\smlbl $\Auth(t)$}};
  \node[evid, below=11mm of r3, minimum width=15mm] (g3)
    {intent\\{\smlbl $\Int(t)$}};
  \node[lbl, below=1mm of g2] {\textbf{governance plane}};

  \draw[eflow] (appr.east) -- (r1.west);
  \draw[eflow] (snap.east) -- (g1.west);

  % the gap
  \draw[bflow] ($(r2.south)+(0,-1mm)$) -- node[lbl, right=1mm, text=black, align=left]
    {unchecked:\\governance drift} ($(g2.north)+(0,1mm)$);
\end{tikzpicture}
}
\caption{Approved intent versus runtime reality. The approval event at $t_0$
freezes a snapshot: configuration, policy version, authorization basis, intent
linkage, and environment assumptions. Afterward, two planes evolve
independently: the runtime plane (observed configuration, running artifacts,
environment) and the governance plane (policy versions, authorization
validity, intent records). Reconcilers continuously compare within the runtime
plane (solid loop); nothing in standard practice compares the runtime plane
against the \emph{evolved} governance plane (dashed gap)---the space in which
governance drift lives.}
\label{fig:planes}
\end{figure}

We fix vocabulary relative to an \emph{approval event}: the moment a change
was legitimized for an environment. \Cref{fig:planes} shows the two planes
that subsequently evolve. All definitions are made precise in the model of
\cref{sec:model}; here we give their content and their separations.

\begin{definition}[Approved State]
\label{def:approved}
The \emph{approved state} of a deployment is an append-only sequence of
immutable snapshots, one per approval event. Each snapshot records: the
approved configuration (manifest content and the
artifact digests it resolves to), the policy version under which it was
evaluated, immutable authorization proof (approvals, exceptions, their
subjects, scopes, execution times, and windows), the intent linkage (the change record and revision the
approval covers), and the declared environment assumptions within which the
risk assessment was made. Evaluation selects the latest applicable covering
snapshot; prior snapshots remain evidence rather than being overwritten.
\end{definition}

\begin{definition}[Configuration Drift]
\label{def:confdrift}
Divergence between the \emph{current desired} configuration and the
\emph{observed} configuration: $\Gobs(t) \ne$ desired$(t)$. This is the
classical notion~\cite{awsconfig,driftctl,nist800128}, decidable by
comparison of present states.
It is an operational component of the umbrella vector; by itself it does not
assert that either state is unauthorized.
\end{definition}

\begin{definition}[Policy Drift]
\label{def:poldrift}
Divergence between the policy version $\pi_B$ pinned in the selected admitted
basis $B_x(t)$ and the policy
context now governing the environment, \emph{as it bears on this deployment}:
$\Pol(t)$ has changed such that the running configuration would not be
admitted under it today, although it was compliant when approved.
\end{definition}

\begin{definition}[Authorization Drift]
\label{def:authdrift}
Invalidation or mismatch of the authorization basis without any configuration
event: a continuing approval or temporary exception lapsed while its effect
persists; an applicable approval was revoked under its declared prospective
or retroactive semantics; or the running artifact's digest is not the subject
any valid approval refers to. A one-shot approval need only have been valid at
execution, provided its integrity-valid proof survives.
\end{definition}

\begin{definition}[Intent Drift]
\label{def:intdrift}
Divergence between the content now driving the runtime and the intent record
the authorization covers: the desired state was changed (e.g., rolled back)
to content whose approval basis is absent or belongs to a different change,
and reconciliation faithfully executes it.
\end{definition}

\begin{definition}[Evidence Drift]
\label{def:evddrift}
Decay of the records connecting the running state to its approval basis: the
chain (revision $\to$ artifact $\to$ deployment $\to$ approval) can no longer
be established from surviving evidence, so governance consistency becomes
\emph{undecidable} rather than false. Evidence drift converts governance
questions from answerable to unanswerable. It is therefore an
\emph{epistemic, transversal} class: unlike the five substantive components,
it does not assert that a governed relation is inconsistent; it records that
one or more such relations can no longer be decided. We keep it inside the
governance-drift umbrella because the operational response---repair evidence
and withhold a polar verdict---is class-specific.
\end{definition}

\begin{definition}[Environment Drift]
\label{def:envdrift}
Violation, outside the managed configuration's scope, of an environment
predicate recorded by the approval's risk assessment: IAM scope broadened,
network exposure altered, an adjacent cloud resource modified out of
band---no desired state exists for the changed property, so no
configuration comparison can even represent the divergence.
\end{definition}

\begin{definition}[Governance Drift]
\label{def:govdrift}
The condition in which the operational state of an environment diverges from
the state, policy context, authorization basis, intent, or environment
assumptions under which it was approved, together with the epistemic condition
in which one or more such relations can no longer be decided. Formally, the
umbrella contains five substantive divergence components
(\cref{def:confdrift,def:poldrift,def:authdrift,def:intdrift,def:envdrift})
and one transversal evidence condition (\cref{def:evddrift}). A deployment is
\emph{governance-consistent} at $t$ when every substantive divergence of
\cref{sec:model} is zero and every required component is decidable. The
substantive components are not mutually exclusive: one incident may
drift several at once (an unapproved rollback drifts intent \emph{and}
leaves the runtime on digests no valid approval covers), which is one reason
the model of \cref{sec:model} keeps them as a vector.
\end{definition}

\subsection{Risk Channels and Severity}
\label{sec:severity}

The five substantive divergence components and the epistemic evidence class
are \emph{not} equivalent in ontology or consequence. A raw count
therefore cannot stand in for severity. Each class opens a different primary
risk channel and implies a different default disposition, as summarized in
\cref{tab:severity}. Configuration drift first threatens operational
integrity; policy drift threatens compliance with the currently governing
control; authorization drift threatens security legitimacy; intent drift
threatens the legitimacy of the change itself; evidence drift threatens audit
and forensic decidability; and environment drift changes the exposure or
assumption boundary around the approved system.

\begin{table*}[t]
\centering
\caption{Proposed Class-Resolved Risk Channels and Illustrative Default
Response Families. These Rows Are Not Empirically Validated or a Scalar
Ranking: Actual Severity Depends on Blast Radius, Exposure Duration,
Reversibility, Consequence, and Evidence Coverage.}
\label{tab:severity}
\footnotesize
\begin{tabular}{@{}p{0.14\textwidth}p{0.23\textwidth}p{0.28\textwidth}p{0.25\textwidth}@{}}
\toprule
Drift class & Primary risk channel & Governance question & Default response family \\
\midrule
Configuration & Operational availability and integrity & Does observed state still match the managed desired state? & Reconcile, contain, or accept an explicit runtime exception \\
Policy & Compliance and control validity & Would the running state satisfy the policy governing it now? & Re-evaluate, migrate under an SLA, or document a current exception \\
Authorization & Security and privilege legitimacy & Is the running subject still covered by valid authority? & Revoke, block, contain, or page the accountable owner \\
Intent & Governance and change legitimacy & Does the deployed revision implement the change that was approved? & Halt propagation, restore the authorized revision, or obtain new approval \\
Evidence & Audit and forensic decidability & Can the approval basis and lineage still be established? & Preserve records, escalate, and report \emph{undecidable} rather than compliant \\
Environment & Assumption and exposure boundary & Do unmanaged surroundings still satisfy the recorded risk assumptions? & Contain exposure and reassess the approval in the changed environment \\
\bottomrule
\end{tabular}
\end{table*}

This mapping is causal, not ordinal: drift class $\rightarrow$ risk channel
$\rightarrow$ response family. Within a class, severity must still be
conditioned on asset criticality, affected subjects, exposure time,
reversibility, and the confidence permitted by surviving evidence. Across
classes, a low-impact authorization mismatch can be less severe than a
configuration fault that stops a safety-critical service. The taxonomy thus
supports class-specific prioritization without pretending that one universal
weight orders every incident.

\begin{figure}[t]
\centering
\resizebox{\columnwidth}{!}{% Drift taxonomy (single column, resized)
\begin{tikzpicture}[node distance=2.8mm and 3mm]
  \node[stage, minimum width=24mm] (conf) at (0,0)
    {\textbf{configuration drift}\\{\smlbl desired vs.\ observed}\\{\smlbl two present configs}\\{\smlbl tier T0n}};
  \node[brok, right=17mm of conf, minimum width=27mm, yshift=9mm] (pol)
    {\textbf{policy drift}\\{\smlbl running state vs.\ $\Pol(t)$}\\{\smlbl tier T1 \; (S3)}};
  \node[brok, below=2.5mm of pol, minimum width=27mm] (auth)
    {\textbf{authorization drift}\\{\smlbl $A_B^{\rm proof}+A_{\rm live}(t)$}\\{\smlbl tier T2/T3 \; (S2, S4, S8)}};
  \node[brok, below=2.5mm of auth, minimum width=27mm] (int)
    {\textbf{intent drift}\\{\smlbl driving content vs.\ approval basis}\\{\smlbl tier T2 \; (S6, jointly authz)}};
  \node[brok, below=2.5mm of int, minimum width=27mm] (evd)
    {\textbf{evidence drift}\\{\smlbl decidability of the above}\\{\smlbl tier T2 \; (S9): verdict $\to \unk$}};
  \node[brok, below=2.5mm of evd, minimum width=27mm] (env)
    {\textbf{environment drift}\\{\smlbl $\varepsilon_t$ vs.\ predicates $\sigma_B$}\\{\smlbl tier T4 \; (S5, S7)}};

  \node[evid, right=8mm of auth.east, minimum width=20mm, yshift=-6mm] (gov)
    {\textbf{governance drift}\\{\smlbl aggregate (Def.~8)}\\{\smlbl $(\GDsub,\ObsMask)$}};

  \draw[flow] (conf.east) to[out=0,in=180] ($(pol.west)+(0,-2mm)$);
  \foreach \n in {pol,auth,int,env}
    \draw[eflow] (\n.east) to[out=0,in=180] (gov.west);
  \draw[bflow] (evd.east) to[out=0,in=180] (gov.west);
  \draw[eflow] (conf.south) to[out=-60,in=200] (gov.south west);
\end{tikzpicture}
}
\caption{The drift taxonomy. Configuration drift (left) is a relation between
two present configurations. Five substantive components relate the present
state to the approved snapshot or evolved context; evidence drift is the
transversal loss of decidability shown by the dotted path. Their umbrella is
governance drift (\cref{def:govdrift}). Annotated
under each class: the evidence tier of \cref{sec:architecture} that first
decides it, as measured in \cref{sec:study}.}
\label{fig:taxonomy}
\end{figure}

\subsection{The Central Separation}
\label{sec:separation}

The taxonomy's load-bearing claim is that \cref{def:govdrift} does not reduce
to \cref{def:confdrift}. We establish it constructively and then in general
form.

\subsubsection{Constructed counterexamples}
\Cref{fig:example} presents the canonical scenario in full: a deployment
whose observed configuration equals its desired configuration at every instant
shown, and which passes from governance-consistent to governance-inconsistent
three separate times---at the policy supersession (policy drift), at the
exception expiry (authorization drift), and at the registry re-tag
(authorization drift at digest granularity). In the fourth panel the desired
state itself is rolled back and the reconciler \emph{eliminates} the
transient configuration drift by converging to content with no valid approval
basis: reconciliation manufactures intent drift while restoring configuration
consistency. Each panel is realizable with standard tooling, and each is one
of the executed scenarios of \cref{sec:study}.

\begin{figure*}[t]
\centering
\resizebox{\textwidth}{!}{% Config-consistent but governance-inconsistent: four-panel timeline (two-column)
\begin{tikzpicture}[x=1mm, y=1mm]
  % timeline base
  \def\W{150}
  \draw[flow] (0,0) -- (\W,0) node[lbl, right]{$t$};
  \node[lbl, left] at (0,0) {events};

  % event marks
  \foreach \x/\lab in {8/{$t_0$: approved\\{\smlbl $\pi_7$, APR-1$\to$\texttt{aaa}, rev-42}},
                       42/{(a) policy\\{\smlbl $\pi_7 \to \pi_8$}},
                       72/{(b) EXC-1\\{\smlbl expires}},
                       102/{(c) re-tag\\{\smlbl 1.4.2$\to$\texttt{bbb}}},
                       130/{(d) git revert\\{\smlbl rev-42$\to$rev-37}}}{
    \draw[black!60] (\x,0) -- (\x,2.5);
    \node[attn, above=0mm] at (\x,3) {\lab};
  }

  % configuration-consistency band
  \node[lbl, left] at (0,-8) {config.\ equality};
  \draw[line width=2.2pt, black] (8,-8) -- (130,-8);
  \draw[line width=2.2pt, black, densely dotted] (130,-8) -- (136,-8);
  \draw[line width=2.2pt, black] (136,-8) -- (\W-4,-8);
  \node[lbl, below] at (70,-9.5) {equal throughout (transient divergence at (d) until reconciler converges)};

  % governance-consistency band
  \node[lbl, left] at (0,-17) {governance cons.};
  \draw[line width=2.2pt, black] (8,-17) -- (42,-17);
  \draw[line width=2.2pt, black, dashed] (42,-17) -- (\W-4,-17);
  \node[lbl, below] at (25,-18.5) {consistent};
  \node[lbl, below] at (60,-18.5) {policy drift};
  \node[lbl, below] at (87,-18.5) {$+$ authz drift};
  \node[lbl, below] at (116,-18.5) {$+$ digest mismatch};
  \node[lbl, below] at (140,-18.5) {$+$ intent drift};

  % reconciler annotation at (d)
  \draw[bflow] (133,-8) to[out=-90,in=90] (136,-15.5);
  \node[lbl, text=black, align=left] at (143,-12.5)
    {reconciler restores\\config equality,\\consummating (d)};
\end{tikzpicture}
}
\caption{Configuration-consistent but governance-inconsistent, in four
panels. Timeline of one deployment: observed configuration (bottom band)
remains equal to desired configuration throughout---configuration drift is
zero at every instant after reconciliation---while governance consistency
(top band) is broken successively by (a)~policy supersession
$\pi_7 \to \pi_8$, (b)~expiry of emergency exception \textsc{exc-1} whose
effect persists, (c)~registry re-tag moving \texttt{svc:1.4.2} to a digest no
approval references, and (d)~a Git rollback that the reconciler faithfully
converges to, eliminating transient configuration drift while stranding the
runtime on content whose approval belongs to a different change. Panels
correspond to Scenarios S3, S2, S4, and S6 of \cref{sec:study}.}
\label{fig:example}
\end{figure*}

\subsubsection{The dependency form}
\begin{proposition}[Configuration equality does not decide governance
consistency]
\label{prop:impossible}
Let $D$ be any detector whose inputs are configuration pairs
$(\mathit{desired}(t), \Gobs(t))$---including their full histories. There
exist worlds $w_1, w_2$ with identical configuration-pair histories in which
the same deployment is governance-consistent in $w_1$ and
governance-inconsistent in $w_2$. Hence no such $D$ decides governance
consistency, for any diffing sophistication.
\end{proposition}
\begin{proof}
Take $w_1$ and $w_2$ to agree on all configurations at all times and differ
only in the governance plane: in $w_2$ the policy version governing the
environment was superseded after $t_0$ by one the running configuration
violates (or: an exception's window elapsed; or: a valid approval in $w_1$ is
absent in $w_2$ for the same content). These differences live in
$\Pol(t), \Auth(t), \Int(t)$, which are not functions of the configuration
pair; both worlds present $D$ with identical inputs, and
\cref{def:govdrift} assigns them different values. $D$ outputs the same answer
in both, so it errs in one.
\end{proof}

The proposition is elementary---a definitional-dependency observation, and
we bill it as such rather than as a deep impossibility theorem. Its force is
architectural: it says deployed drift detectors are not approximations of
governance-drift detectors but detectors \emph{of a different relation}, and
no incremental improvement of configuration diffing closes the gap
\emph{as a decision procedure}. Two bounds on what it claims: it rules out
\emph{deciding}, not statistical \emph{hinting}---in real estates,
governance-plane events often correlate with configuration-plane events
(policy revisions ship alongside config pushes; rollbacks are visible in Git
history), so configuration histories may carry probabilistic signal about
governance state; how much is an empirical question for the laboratory, which
our closed world deliberately does not answer. And it says nothing against
comparison as a \emph{mechanism}: the governance tiers of
\cref{sec:architecture} are themselves comparisons---against recorded
bases, validity clocks, and coverage sets---which is precisely the point:
the input vocabulary, not the mechanism, is what must widen.

\subsubsection{What configuration drift remains}
The separation does not demote configuration drift: unauthorized manual
change (\cref{sec:study}, S1) is both configuration and governance drift, and
configuration comparison is its right detector. The taxonomy's point is
containment, not replacement: configuration drift is one of five substantive
components under the six-class umbrella, and the only one decidable in the
configuration vocabulary---which the study confirms
empirically as a detectability staircase whose first step is exactly that one
class.
